\documentclass[11pt,a4paper]{article}
\usepackage[T1]{fontenc}
\usepackage[utf8]{inputenc}
\usepackage{lmodern}
\usepackage[margin=2.5cm]{geometry}
\usepackage{booktabs}
\usepackage{tabularx}
\usepackage{longtable}
\usepackage{enumitem}
\usepackage{listings}
\usepackage{xcolor}
\usepackage[hidelinks]{hyperref}

\lstset{
  basicstyle=\ttfamily\small,
  columns=fullflexible,
  keepspaces=true,
  frame=single,
  rulecolor=\color{gray!50},
  backgroundcolor=\color{gray!8},
  breaklines=true,
  upquote=true,
  showstringspaces=false,
  extendedchars=true,
  inputencoding=utf8,
  literate={é}{{\'e}}1 {è}{{\`e}}1 {à}{{\`a}}1 {ç}{{\c c}}1 {É}{{\'E}}1,
  xleftmargin=0.5em,
  xrightmargin=0.5em
}
\setlength{\emergencystretch}{3em}
\newcommand{\kbd}[1]{\texttt{#1}}
\newcolumntype{L}{>{\raggedright\arraybackslash}X}

\title{LeRobot (Free Pascal) User Manual}
\author{}
\date{9 November 2026}

\begin{document}
\maketitle

LeRobot is a drawing robot you steer on screen: either live, with the mouse
and keyboard (Dessiner), or by writing a short Pascal program of robot
commands (such as \kbd{Dessin\_du\_Robot.p}). This version compiles the
Turbo Pascal programs of the 1991 book \emph{Premières leçons de
programmation en Turbo Pascal} (Macintosh) with Free Pascal, unchanged, and
runs them in the robot window of the OCaml and Python versions, with the
same behaviour. It lives in
the folder \kbd{LeRobot\_FreePascal\_FR} of the \kbd{LeRobot} folder, which can
be anywhere (commands below are typed in a Terminal window opened in the
\kbd{LeRobot} folder); the installation
is described in the \emph{LeRobot (Free Pascal) Installation Manual}
(\kbd{LeRobot\_Install.pdf}).

\tableofcontents

%======================================================================
\section{Getting started}

You need Free Pascal (\kbd{fpc}), a Python linked to Tk 8.6 or later (for
example the one from python.org), and the folders \kbd{LeRobot\_FreePascal\_FR}
and \kbd{LeRobot\_Python\_FR} side by side (see the installation manual). There
is no XQuartz: the robot window is an ordinary macOS window.

There are two ways to start a program:
\begin{itemize}
\item \textbf{The application} \kbd{LeRobotFreePascal.app} (in
  \kbd{LeRobot\_FreePascal\_FR}): double-click it to run Dessiner; drop a
  \kbd{.p} file on its icon to compile and run that program.
\item \textbf{The script} \kbd{robot.sh}, in a Terminal window opened in
  \kbd{LeRobot\_FreePascal\_FR}: \kbd{./robot.sh path/to/Program.p}. The
  program can be in any folder.
\end{itemize}
Both compile the program with Free Pascal (the compiler messages appear
only if there is an error), then run it. The compiled programs are kept in
the folder \kbd{\_build}, which can be deleted at any time.

\paragraph{How it works.} The Macintosh of 1991 drew with QuickDraw, which
does not exist any more. The units \kbd{Robot1}, \kbd{Robot2} and
\kbd{Robot3} of this version send each command of the program (\kbd{av},
\kbd{lg(12)}, \dots) to the robot window of the Python version, which draws
it and answers. The program itself is an ordinary compiled Pascal program:
\kbd{write}, \kbd{writeln} and \kbd{readln} use the Terminal window.

The robot window is $800 \times 700$ points. The large white area is where
the robot draws; the strip along the bottom is the message bar.

%======================================================================
\section{Part 1 --- Interactive drawing (Dessiner)}

In Dessiner you give the robot commands with single keys and move it with
the mouse or the arrow keys; press Enter when you are done, then \kbd{y} to
save the drawing as a Pascal program.

\subsection{Launch}
Double-click \kbd{LeRobotFreePascal.app}, or:
\begin{lstlisting}[language=bash]
cd LeRobot_FreePascal_FR
./robot.sh livre/Didactitiels_et_jeux/Dessiner.p
\end{lstlisting}
\kbd{Dessiner.p} is the program of 1991: it calls the procedure
\kbd{interaction} of the robot with the drawing and saving options. The
robot appears in the centre, pointing north. The window comes to the front
with the keyboard focus; if keys do nothing, click in the window.

\subsection{The window}
\begin{itemize}
\item \textbf{Drawing area} --- the large white area.
\item \textbf{Message bar} (bottom) --- a reminder of the keys, ending with
  \kbd{* :cc w :peindre\ \ Entree :fin}.
\item \textbf{Colour box} (small robot near the right of the bar) --- shows
  the current pen colour.
\item \textbf{Enter icon} (far right) --- the Enter symbol; it blinks when
  the robot waits for a key.
\end{itemize}

\subsection{Moving the robot with the mouse or the arrow keys}
Click on a grid point (or move the pointer with the button held down) and
the robot jumps to the nearest grid point, without drawing. The arrow keys
move it one grid step left, right, up or down, also without drawing. When
the pointer moves with the button released, the robot does not move.

If the robot would go past the edge of the drawing area, it stays where it
is and the terminal prints \kbd{** SORTIE INTERDITE ** translater}
(``forbidden exit''), with a beep-beep. This is only a warning.

\subsection{Keyboard commands}
The robot draws only when its pen is down (\kbd{b}; it is down at the
start). One step is one grid square.

\begin{center}
\begin{tabularx}{\linewidth}{@{}l l L@{}}
\toprule
\textbf{Key} & \textbf{Command} & \textbf{What it does} \\
\midrule
\kbd{b} & bc & Pen down: the robot draws as it moves \\
\kbd{l} & lc & Pen up: the robot moves without drawing \\
\kbd{a} & av & Move forward one step \\
\kbd{)} & vd & Curved turn to the right (draws an arc) \\
\kbd{(} & vg & Curved turn to the left (draws an arc) \\
\kbd{+} & pqtd & Quarter turn right ($90^\circ$) \\
\kbd{-} & pqtg & Quarter turn left ($90^\circ$) \\
\kbd{d} & phtd & Eighth turn right ($45^\circ$) \\
\kbd{g} & phtg & Eighth turn left ($45^\circ$) \\
\kbd{c} & dc & Draw a small cross (+) \\
\kbd{p} & dp & Draw a point \\
\kbd{.} & dg & Draw the grid points \\
\kbd{*} & cc & Next pen colour (shown in the colour box) \\
\kbd{w} & peindre & Paint the closed area behind the robot with the pen colour \\
Arrow keys & translater & Move one step left, right, up or down, without drawing \\
Backspace & --- & Undo the last command \\
Enter & --- & Finish the drawing \\
\bottomrule
\end{tabularx}
\end{center}

Colours cycle in this order: white, yellow, blue, violet, green, red,
indigo, black. To paint, put the robot so that its back is inside a closed
area, then press \kbd{w}.

\subsection{Finishing and saving}
\begin{enumerate}
\item Press \textbf{Enter}. The bar asks:
  \emph{Voulez-vous sauver dans \dots/Dessin\_du\_Robot.p ? (y/n)}
\item Press \kbd{y} to save. Every command you used, including mouse and
  arrow moves (written as \kbd{translater(dx, dy)}) and colour changes
  (written as \kbd{cc(n)} with the colour name in a comment), is written to
  \kbd{Dessin\_du\_Robot.p} as a Pascal program (see Part~2), and the
  window closes. The Terminal prints the command to replay it.
\item Press any other key to close without saving.
\end{enumerate}
Saving always overwrites \kbd{Dessin\_du\_Robot.p}, in the
\kbd{LeRobot\_FreePascal\_FR} folder. To keep an earlier drawing, copy that
file under another name first.

%======================================================================
\section{Part 2 --- Running a robot program}

A robot program is a Pascal program that uses the unit \kbd{Robot1} and
lists robot commands in order. It draws the picture by itself, then waits
so you can replay it at another speed or close it with Esc.

\subsection{Run Dessin\_du\_Robot.p}
\begin{lstlisting}[language=bash]
cd LeRobot_FreePascal_FR && ./robot.sh Dessin_du_Robot.p
\end{lstlisting}
or drop \kbd{Dessin\_du\_Robot.p} on \kbd{LeRobotFreePascal.app}. A saved
drawing looks like this:
\begin{lstlisting}[language=Pascal]
{$mode macpas}
program Dessin_du_Robot;
 uses Robot1, Robot2;
 { Mon joli dessin, engendre par Dessiner }
begin
 ce;
 av; av; pqtd; av; av; cc(2); { bleu }
 av; translater(5, 3); av; vg; av; st;
end.
\end{lstlisting}

\subsection{Write your own program}
Create a file anywhere, for example \kbd{Maison.p}:
\begin{lstlisting}[language=Pascal]
program Maison;
 uses Robot1;
 { Dessiner un carre }
begin
 av; av; pqtd; av; av; pqtd;
 av; av; pqtd; av; av; pqtd;
 st;
end.
\end{lstlisting}
Run it with \kbd{./robot.sh path/to/Maison.p} from the
\kbd{LeRobot\_FreePascal\_FR} folder, or drop it on \kbd{LeRobotFreePascal.app}.

Every program follows the same pattern:
\begin{itemize}
\item \kbd{uses Robot1;} after the \kbd{program} line. All the commands are
  in \kbd{Robot1}; \kbd{Robot2} and \kbd{Robot3} (sounds, texts, input in
  1991) are empty and may be listed too, as in the programs of the book.
  The window opens at the first command, with the robot in the centre,
  pointing north, pen down, in black, with a grid of 12 points.
\item The robot starts with its pen down, so every move draws: call
  \kbd{lc} to move without drawing, and \kbd{bc} to draw again.
\item \kbd{st} last: shows the speed screen and keeps the window open.
  Without it, the bar says \emph{Fin du programme : taper une touche ou
  cliquer} and the window closes at the next key or click.
\item Commands without arguments are written without parentheses
  (\kbd{av}, not \kbd{av()}); commands are separated by \kbd{;}. Pascal
  does not distinguish upper and lower case: \kbd{AV}, \kbd{Av} and
  \kbd{av} are the same command.
\item The program may be written in UTF-8 (accents allowed in strings and
  comments) or come from the Macintosh of 1991 (MacRoman encoding, Mac
  units such as \kbd{Memtypes}, \kbd{Quickdraw}): \kbd{robot.sh} converts
  the latter first, into \kbd{\_build/convertis}, without changing the
  original. The line numbers of the error messages stay those of the
  original.
\end{itemize}

\subsection{Errors}
\begin{itemize}
\item \textbf{Compilation errors} are found by Free Pascal before the
  program runs; the window does not open. The message gives the line and
  column:
\begin{lstlisting}
Maison.p(5,2) Error: Identifier not found "avv"
Maison.p(5,2) Error: Wrong number of parameters specified for call to "av"
\end{lstlisting}
\item \textbf{Leaving the window}: when the robot would go past the edge,
  the program stops with a beep-beep and the Terminal prints
  \kbd{Erreur : le robot sortirait du cadre (av)}; the window stays open
  to show the drawing made so far.
\end{itemize}

\subsection{Programs of the book}
The folder \kbd{livre} holds the 407 programs of the book, one subfolder per
chapter (the same names as in the OCaml and Python versions):
\begin{lstlisting}[language=bash]
./robot.sh livre/20_grille_polaire/LimaconDePascal.p
./verifier.sh                         # all programs, one after the other
./verifier.sh 02_interaction          # one chapter
./verifier.sh '' FriseDeTriangles     # resume at a program
\end{lstlisting}
With \kbd{verifier.sh}, answer in the robot window when a drawing is
finished: Enter = correct (next program), \kbd{r} = run it again, \kbd{q}
or Esc = stop (the command to resume is printed), any other key starts a
remark, ended by Enter, noted in \kbd{verification.txt}. Compilation errors
are noted there too (\kbd{Cocotte-errone.p} has one on purpose, as in the
book). A program that never ends by itself (\kbd{Boucle1}, \kbd{Boucle2},
\kbd{MontreDigitale}) is stopped with Esc. Programs that ask questions
(\kbd{readln}) ask them in the Terminal: click in the Terminal window to
answer. \kbd{LEROBOT\_MUET=1} turns the sound off.

The programs are the Macintosh programs of 1991, converted once by \kbd{outils/convertir.py} (UTF-8, units of the robot, directive
\kbd{\{\$mode macpas\}} for the Macintosh dialect of Pascal); their text is
unchanged. Some of them were changed later in the OCaml and Python editions
of the book, so a few drawings differ from those versions.

\subsection{Replaying at another speed}\label{sec:replay}
When the program reaches \kbd{st}, the bar shows:
\begin{lstlisting}
+ : plus vite   - : moins vite   vitesse : N     autre touche : redessiner   Echap : quitter
\end{lstlisting}
\begin{itemize}
\item \kbd{+} / \kbd{-} --- change the speed from 0 to 10 (default 10, the
  fastest). Only the number changes.
\item Any other key, for example the space bar --- clear the window and
  redraw at the chosen speed; the speed screen comes back when it is done.
\item Esc --- close the window.
\end{itemize}
At speed 1 each command takes about 0.1\,s. Speed 0 is step-by-step mode:
the robot does one command each time you press a key. While the robot
draws, Esc (or \kbd{\S} or \kbd{\#}) stops the program.

\subsection{Command reference}
One step is one grid square. Commands are in French; the table gives the
English meaning. The numbers in brackets are parameters (\kbd{integer} or
\kbd{real}); \kbd{var} parameters receive a result.

\begin{longtable}{@{}p{0.42\linewidth} p{0.53\linewidth}@{}}
\toprule
\textbf{Command} & \textbf{Meaning} \\
\midrule
\endhead
\bottomrule
\endfoot
\kbd{av} & Move forward one step \\
\kbd{avf(n)} & Move forward $n$ steps (backwards if $n < 0$) \\
\kbd{vd} / \kbd{vg} & Curved turn right / left (draws an arc) \\
\kbd{phtd} / \kbd{phtg} & Turn $45^\circ$ right / left \\
\kbd{pqtd} / \kbd{pqtg} & Turn $90^\circ$ right / left \\
\kbd{p3htd} / \kbd{p3htg} & Turn $135^\circ$ right / left \\
\kbd{pdt} & Turn around ($180^\circ$) \\
\kbd{pn} \kbd{pne} \kbd{pe} \kbd{pse} \kbd{ps} \kbd{psw} \kbd{pw} \kbd{pnw} & Face N, NE, E, SE, S, SW, W, NW \\
\kbd{pivoter(n)} & Face direction $n$ (\kbd{nord} = 0, \kbd{nordEst} = 1 \dots\ \kbd{nordOuest} = 7) \\
\kbd{rt(a)} & Turn by $a$ degrees (anticlockwise), keeping the step length \\
\kbd{bc} / \kbd{lc} & Pen down / pen up \\
\kbd{dp} / \kbd{dc} & Draw a point (a square of side \kbd{ec}) / a small cross \\
\kbd{cc(n)} & Pen colour $n$ (\kbd{blanc} = 0, \kbd{jaune}, \kbd{bleu}, \kbd{violet}, \kbd{vert}, \kbd{rouge}, \kbd{indigo}, \kbd{noir} = 7) \\
\kbd{ec(n)} & Pen thickness $n$ (a square pen, as on the Macintosh) \\
\kbd{cf(n)} & Clear the window and fill it with background colour $n$ \\
\kbd{peindre} & Paint the closed area behind the robot with the pen colour \\
\kbd{translater(dx, dy)} & Move $dx$ steps across and $dy$ steps up, without drawing \\
\kbd{tracer(x, y)} & Go to window pixel $(x, y)$, drawing if the pen is down \\
\kbd{deplacer(x, y)} & Go to grid point $(x, y)$, counted from the centre \\
\kbd{tl(l)} & Move forward $l$ points, the grid following the robot \\
\kbd{ce} & Go to the centre \\
\kbd{eh} \kbd{eb} \kbd{ad} \kbd{ag} & Go to the top, bottom, right or left edge \\
\kbd{dg} & Draw the grid points \\
\kbd{lg(x)} / \kbd{lgX(x)} / \kbd{lgY(x)} & Set the step size (both directions / across / up) \\
\kbd{lgRT(r, t)} & Step of length $r$ at angle $t$ (polar grid) \\
\kbd{montrerRobot} / \kbd{cacherRobot} & Show / hide the robot \\
\kbd{autoriserSorties} / \kbd{interdireSorties} & Allow / forbid moving past the window edge \\
\kbd{vt(n)} & Set the speed (0 to 10) \\
\kbd{pp} / \kbd{ex} & Step-by-step mode / full speed \\
\kbd{message('text')} & Write a message in the bottom bar \\
\kbd{ecrire('text')} & Write a text at the robot's position \\
\kbd{entierEnChaine(n)}, \kbd{reelEnChaine(x)} & A number as a string, for \kbd{message} and \kbd{ecrire} \\
\kbd{lireCar(c)} & Wait for a key in the robot window; its character is put in \kbd{c} (a \kbd{var} parameter of type \kbd{char}) \\
\kbd{interaction(x, y, c, \dots)} & Let the user move the robot (see \kbd{Dessiner.p}) \\
\kbd{note(h, d)}, \kbd{silence(d)}, \kbd{bip} & Play a note, a rest, a beep \\
\kbd{delai(t)} & Wait $t$ seconds \\
\kbd{heure(h, m, s)}, \kbd{date(j, m, a)} & The time, the date (\kbd{var} parameters) \\
\kbd{write}, \kbd{writeln}, \kbd{readln} & Write and read in the Terminal (standard Pascal) \\
\kbd{st} & End: show the speed screen \\
\end{longtable}

Test functions such as \kbd{tb} (would the next step leave the window?)
and the values \kbd{valX}, \kbd{valLgX}, \kbd{valOrientation}, \dots\ are
listed with all the other commands, constants (colours, directions, notes)
and their parameter types in the interface of
\kbd{unites/Robot1.pas}. The architecture and the conversion of the
programs are described in the \kbd{README.md} file of
\kbd{LeRobot\_FreePascal\_FR}.

%======================================================================
\section{Troubleshooting}

\begin{longtable}{@{}p{0.27\linewidth} p{0.33\linewidth} p{0.34\linewidth}@{}}
\toprule
\textbf{Symptom} & \textbf{Cause} & \textbf{Fix} \\
\midrule
\endhead
\bottomrule
\endfoot
\emph{Prog.p(5,2) Error: Identifier not found "avv"} & A command misspelled (or a missing \kbd{uses Robot1;}) & Correct line 5, column 2 \\
\addlinespace
\emph{Wrong number of parameters specified for call to "av"} & \kbd{av()} or \kbd{av(3)} written for \kbd{av} & Write \kbd{av} without parentheses, or \kbd{avf(3)} \\
\addlinespace
\emph{Erreur : le robot sortirait du cadre (av)} & The robot was asked to move past the window edge & Change the program, or call \kbd{autoriserSorties} \\
\addlinespace
\kbd{** SORTIE INTERDITE ** translater} in the terminal & In Dessiner, a move past the window edge & Only a warning \\
\addlinespace
The program never ends & An endless loop (\kbd{while true}, \kbd{repeat \dots\ until false}, endless recursion) & Press Esc in the robot window \\
\addlinespace
The program waits and the window shows nothing new & It asks a question with \kbd{readln} in the Terminal & Click in the Terminal window and answer \\
\addlinespace
Keys do nothing in the window & The window does not have keyboard focus & Click in the window \\
\addlinespace
Digits are not accepted & French AZERTY keyboard: the top row gives \& é " \dots\ without Shift & For \kbd{lireCar}, use Shift or the numeric keypad \\
\addlinespace
\emph{Ce programme doit etre lance par robot.sh} & The compiled program was started by itself & Use \kbd{./robot.sh Prog.p} or \kbd{LeRobotFreePascal.app} \\
\addlinespace
A saved drawing was lost & Saving from Dessiner always overwrites \kbd{Dessin\_du\_Robot.p} & Copy it under a new name before saving again \\
\addlinespace
Text font cannot be changed & \kbd{policeTexte} and \kbd{styleTexte} have no effect & None; texts use Helvetica 13 \\
\end{longtable}

\end{document}
